\clearpage
\fancyhead[L]{\small VON LINDENAU. STERNWARTE SEEBERG.}
\hypertarget{native-503-lower}{}\label{native-503-lower}
\newcommand{\SeebergResidualMark}{\hbox{\vrule width0pt height3.272222bp depth0pt\special{pdf:bcontent}\special{pdf:code q 0 g 2.927778 1.377778 m 2.927778 1.205556 l 2.755556 1.033333 l 2.755556 0.172222 l 2.927778 0.000000 l 3.100000 0.172222 l 3.100000 1.377778 l h 0.172222 3.272222 m 0.000000 3.100000 l 0.000000 2.927778 l 0.172222 2.755556 l 1.205556 2.755556 l 1.377778 2.927778 l 2.066667 2.927778 l 2.238889 2.755556 l 2.755556 2.755556 l 2.927778 2.583333 l 2.927778 2.411111 l 3.100000 2.411111 l 3.100000 2.583333 l 3.444444 2.927778 l 6.027778 2.927778 l 6.200000 3.100000 l 6.027778 3.272222 l 3.961111 3.272222 l 3.788889 3.100000 l 3.616667 3.100000 l 3.444444 3.272222 l 2.583333 3.272222 l 2.411111 3.100000 l 0.516667 3.100000 l 0.344444 3.272222 l h f* Q}\special{pdf:econtent}\kern6.200000bp}}
\newcommand{\SeebergMoon}{\text{\mdseries\AstroSymbols ☽}}

\begin{center}
\small Göttingische gelehrte Anzeigen. Stück 106. S. 1049..1052. 2. Juli 1808.\\[3pt]
\rule{.72\linewidth}{.3pt}
\end{center}

Herr \textsc{von Lindenau} hat uns kürzlich in einem eigenen Aufsatz einen zuverlässigen Bericht über den gegenwärtigen Zustand der berühmten Sternwarte \emph{Seeberg} zugesandt. Ein Auszug daraus wird den Freunden der Himmelskunde umso willkommener sein, als seit einiger Zeit manche unbegründeten Nachrichten darüber aus trüben Quellen an die Öffentlichkeit gelangt waren. Die Kriegsunruhen und einige notwendig gewordene Reparaturen am Dach der Sternwarte hatten 1806 den Abbau der Instrumente erforderlich gemacht. Verschiedene Umstände, die hier nicht aufzuzählen sind, ließen es erst zu Beginn dieses Jahres zu, die nötigen Vorkehrungen für die Wiederaufstellung zu treffen. Herr \textsc{von Lindenau}, der in Abwesenheit von Herrn \textsc{von Zach} schon einmal die Aufsicht über die Sternwarte geführt hatte, wurde von Seiner Durchlaucht, dem Herzog von Sachsen-Gotha, mit dieser Aufgabe betraut und nahm am 9. April seine Wohnung auf der Sternwarte. Wegen des ungünstigen Wetters konnte man erst am 20. April mit der Wiederaufstellung der Instrumente beginnen. Danach begünstigten jedoch mehrere klare Nächte die Arbeit so sehr, dass Passageninstrument, Quadrant und Regulator Ende April für die Beobachtungen bereit waren. Wie vollkommen Herrn \textsc{von Lindenau} die Justierung der Instrumente sogleich gelang, zeigt eine schöne Reihe zahlreicher Beobachtungen, die er vom 29. April bis 9. Mai anstellte und uns mitteilte. Wegen des beschränkten Platzes in diesen Blättern können wir hier nur einige davon wiedergeben.

\begin{center}
\emph{Beobachtungen der Sonne.}

\smallskip
{\small\setlength{\tabcolsep}{3pt}
\begin{tabular}{r|rrr|rrr|r}
1808 & \multicolumn{3}{c|}{\shortstack{Mittlere Zeit\\in Seeberg}} & \multicolumn{3}{c|}{\shortstack{Rektaszension\\der Sonne}} & \shortstack{Fehler der\\Tafeln} \\
\hline
April 30 & $23^{\mathrm{u}}$ & $57^{\mathrm{m}}$ & $4^{\mathrm{s}}16$ & $37^{\circ}$ & $29'$ & $4''2$ & $+3''8$ \\
Mai 3 & 23 & 56 & 42.06 & 40 & 20 & 57.4 & $+1.0$ \\
5 & 23 & 56 & 29.87 & 42 & 16 & 9.9 & $+2.5$ \\
6 & 23 & 56 & 24.89 & 43 & 14 & 3.3 & $-2.0$ \\
7 & 23 & 56 & 20.06 & 44 & 11 & 59.0 & $-0.5$ \\
9 & 23 & 56 & 11.73 & 46 & 8 & 15.5 & $\SeebergResidualMark\,2.6$ \\
\hline
\multicolumn{7}{r|}{Fehler der Tafeln im Mittel} & $+1.8$ \\
\end{tabular}}
\end{center}

\clearpage
\SourceLeaf{504}
\begin{center}
\emph{Beobachtungen des Mondes.}

\smallskip
{\small\setlength{\tabcolsep}{3pt}
\begin{tabular}{r|rrr|rrr|r}
1808 & \multicolumn{3}{c|}{\shortstack{Mittlere Zeit\\in Seeberg}} & \multicolumn{3}{c|}{\shortstack{Rektaszension des\\westl. $\SeebergMoon$-Randes}} & \shortstack{Fehler der\\Tafeln.} \\
\hline
Mai 2 & $5^{\mathrm{u}}$ & $20^{\mathrm{m}}$ & $37^{\mathrm{s}}63$ & $120^{\circ}$ & $34'$ & $0''3$ & $+2''5$ \\
3 & 6 & 10 & 48.31 & 134 & 7 & 50.4 & $+0.6$ \\
4 & 7 & 0 & 59.58 & 147 & 41 & 52.8 & $+5.7$ \\
5 & 7 & 51 & 44.20 & 161 & 24 & 15.9 & $+6.2$ \\
6 & 8 & 43 & 44.35 & 175 & 25 & 30.6 & $+3.1$ \\
7 & 9 & 37 & 44.89 & 189 & 57 & 4.3 & $+2.3$ \\
8 & 10 & 13 & 19.15 & 205 & 7 & 5.8 & $+6.7$ \\
\hline
\multicolumn{7}{r|}{Mittlerer Fehler der Tafeln} & $+3''9$ \\
\end{tabular}}
\end{center}

Die Tafeln, mit denen die Beobachtungen verglichen wurden, sind die neuen Sonnentafeln von Herrn \textsc{von Zach} und die Mondtafeln von \textsc{Bürg}.

Außerdem teilte uns Herr \textsc{von Lindenau} noch 8 beobachtete Positionen des \emph{Uranus} mit, 11 Positionen des \emph{Saturn}, 6 Positionen des \emph{Jupiter}, 9 Positionen der Venus und 8 Positionen des \emph{Merkur}. Wir begnügen uns damit, hier nur die Endergebnisse anzuführen, die Herr \textsc{von Lindenau} daraus gewonnen hat. Der mittlere Fehler der Uranustafeln von \textsc{de Lambre} betrug für diese Epoche $-0''8$ in der Länge und $-16''$ in der Breite. Der mittlere Fehler der Saturntafeln betrug in der Länge $+9''1$ und in der Breite $-4''3$. Die Opposition des Saturn fand 1808 am 9. Mai um $8^{\mathrm{u}}25^{\mathrm{m}}47^{\mathrm{s}}2$ mittlerer Zeit in Seeberg statt. Der mittlere Fehler der Jupitertafeln betrug in der Länge $-1''6$ und in der Breite $-1''9$. Die Beobachtungen der Venus und des Merkur hat Herr \textsc{von Lindenau} nicht verglichen, weil er die Tafeln dieser Planeten dem heutigen Stand der Astronomie nicht mehr angemessen fand. Seit einem Jahr sammelt er selbst Material für neue Venustafeln. Dafür hat er bereits mehr als 200 gute Beobachtungen zusammengetragen. Damit und mithilfe der von \textsc{la Place} entwickelten Störungsgleichungen hofft er bald etwas Vollkommeneres liefern zu können als die bisherigen Tafeln.

Als Hauptaufgabe seiner praktischen Arbeit während seines Aufenthalts auf der Seeberger Sternwarte hat sich Herr \textsc{von Lindenau} vorgenommen, die Refraktionen und die jährliche Parallaxe der Fixsterne in Rektaszension zu bestimmen. Wenn die Annahme begründet ist, dass $\alpha$ Lyrae eine Parallaxe in Deklination von $4\ldots5''$ zeigt, so muss die Parallaxe in Rektaszension $7''$ betragen. Der Unterschied zwischen dem positiven und dem negativen Maximum steigt also auf fast eine Zeitsekunde. Vergleicht man einen solchen Stern mit mehreren kleinen Sternen in seiner Nähe, so darf man hoffen, dass eine solche Differenz dem vortrefflichen 8 Fuß langen Passageninstrument von \textsc{Ramsden} nicht entgehen wird. Freilich wird der Vergleich mit kleinen Sternen im positiven und negativen Maximum beträchtlichen \emph{praktischen} Schwierigkeiten unterliegen.

Herr \textsc{von Lindenau} hat mit den Refraktionsbeobachtungen bereits begonnen. Obwohl es noch zu früh ist, daraus ein bestimmtes Ergebnis zu ziehen, ergeben sie doch alle bei einem Barometerstand von 28 Zoll und einer Temperatur von $10^{\circ}$ Réaumur eine Horizontalrefraktion von $37\ldots40'$. (Wir bemerken bei dieser Gelegenheit, dass auch Herr Inspektor \textsc{Bessel} aus einer sehr sorgfältigen Untersuchung der Beobachtungen von \textsc{Bradley} geschlossen hat, dass der übliche Wert für die Horizontalrefraktion beträchtlich vergrößert werden muss). — Insgesamt wird Herr \textsc{von Lindenau} den Astronomen nach Ablauf eines Jahres die Ergebnisse seines Aufenthalts auf der Seeberger Sternwarte ausführlich vorlegen.

Das ganze Gebäude der eigentlichen Sternwarte ist jetzt vollständig wiederhergestellt. Am Wiederaufbau des Wohngebäudes wird gearbeitet, sodass es hoffentlich ebenfalls noch im Laufe dieses Sommers fertig wird.
